% Part: incompleteness
% Chapter: arithmetization-syntax
% Section: proofs-in-ax

\documentclass[../../../include/open-logic-section]{subfiles}

\begin{document}

\olfileid{inc}{art}{pax}
\olsection{公理化推导}

\begin{explain}
为了将公理推导算术化，我们必须把推导表示为数。
由于推导仅仅是公式的序列，显而易见的做法是将每个推导编码为该推导中各公式编码所构成序列的码。
\end{explain}

\begin{defn}
若 $\delta$ 是一个由公式 $!A_1$, \dots.~$!A_n$ 组成的公理推导，则 $\Gn{\delta}$ 是
\[
\tuple{\Gn{!A_1}, \dots, \Gn{!A_n}}.
\]
\end{defn}

\begin{ex}
  考虑下面这个非常简单的推导：
  \begin{derivation}
  1. & $!B \lif (!B \lor !A)$ \\
  2. & $(!B \lif (!B \lor !A)) \lif (!A  \lif (!B \lif (!B \lor !A)))$\\
  3. & $!A  \lif (!B \lif (!B \lor !A))$
  \end{derivation}
  此推导的Gödel数将是
  \begin{align*}
  \openTuple\,
    & \Gn{!B \lif (!B \lor !A)}, \\
    &\Gn{(!B \lif (!B \lor !A)) \lif (!A  \lif (!B \lif (!B \lor !A)))},\\
    & \Gn{!A  \lif (!B \lif (!B \lor !A))} \,\closeTuple.
  \end{align*}
\end{ex}

\begin{explain}
确定了推导的表示方法后，我们还必须证明能够原始递归地操作这些推导，
并原始递归地表达其基本性质与关系。
有些操作很简单：例如，给定一个推导的Gödel数~$d$，
$(d)_{\len{d}-1}$ 给出其末尾公式的Gödel数。
有些则困难得多。我们至少会概述如何做到这一点。目标是证明关系
“$\delta$ 是从~$\Gamma$ 到~$!A$ 的推导”在 $\delta$ 和 $!A$ 的Gödel数上是原始递归的。
\end{explain}

\begin{prop}
\ollabel{prop:followsby}
下列关系是原始递归的：
\begin{enumerate}
\item $!A$ 是一条公理。
\item $\delta$ 中的第 $i$ 行由肯定前件证成。
\item $\delta$ 中的第 $i$ 行由量词规则 \QR 证成。
\item $\delta$ 是一个正确的推导。
\end{enumerate}
\end{prop}

\begin{proof}
我们必须证明，公式与推导的Gödel数之间的相应关系是原始递归的。
\begin{enumerate}
\item 我们有一个给定的公理模式列表，$!A$ 是一条公理，如果它符合其中某个模式的形式。
  由于模式列表有限，只需证明，对于每个公理模式，
  我们都能原始递归地测试 $!A$ 是否具有该形式。
  例如，考虑公理模式
  \[
  !B \lif (!C \lif !B).
  \]
  $!A$ 是该公理模式的一个实例，如果存在公式 $!B$ 和 $!C$，
  使得将 `$($'、$!B$、`$\lif$'、`$($'、$!C$、`$\lif$'、$!B$ 与 `$))$' 拼接起来能得到 $!A$。
  我们能够对 $!A$ 的Gödel数 $n$ 测试相应的性质，因为序列的拼接是原始递归的，
  且 $!B$ 和 $!C$ 的Gödel数必定小于 $!A$ 的Gödel数——
  因为当该关系成立时，$!B$ 和 $!C$ 都是 $!A$ 的子公式。
  因此，我们可以定义：
  \begin{multline*}
  \fn{IsAx}_{!B \lif (!C \lif !B)}(n) \defiff \bexists{b< n}{
    \bexists{c < n}{(\fn{Sent}(b) \land \fn{Sent}(c) \land {}}}\\ n =
  \Gn{(} \concat b \concat \Gn{\lif} \concat \Gn{(} \concat c \concat
  \Gn{\lif} \concat b \concat \Gn{))}).
  \end{multline*}
  如果我们对每个公理模式都有这样一个定义，那么它们的析取就定义了性质
  $\fn{IsAx}(n)$，即“$n$~是某条公理的Gödel数”。

\item $\delta$ 中的第 $i$ 行由肯定前件证成，当且仅当存在行 $j$ 和 $k < i$，
  使得第 $j$ 行是某个公式 $!A$，第 $k$ 行是 $!A \lif !B$，
  而第 $i$ 行是 $!B$。
  \begin{multline*}
    \fn{MP}(d, i) \defiff \bexists{j < i}{\bexists{k < i}{}}\\
    (d)_k = \Gn{(} \concat (d)_j
      \concat \Gn{\lif} \concat (d)_i \concat \Gn{)}
  \end{multline*}
  因为有界量化、拼接和 $=$ 是原始递归的，这定义了一个原始递归关系。

\item $\delta$ 中的一行由 \QR{} 证成，如果它具有形式 $!B \lif \lforall[x][!A(x)]$，
  且前面有一行是 $!B \lif !A(c)$（对于某个常项符号 $c$），并且 $c$ 不在 $!B$ 中出现。
  这成立当且仅当：
  \begin{enumerate}
  \item 存在一个语句 $!B$，且
  \item 存在一个仅含单自由变元 $x$ 的公式 $!A(x)$，使得
  \item 第 $i$ 行包含 $!B \lif \lforall[x][!A(x)]$
  \item 某个第 $j < i$ 行包含 $!B \lif \Subst{!A}{c}{x}$（对于某个常项 $c$），
  \item $c$ 不在 $!B$ 中出现。
  \end{enumerate}
  所有这些都可以原始递归地测试，因为 $!B$、$!A(x)$ 和 $x$ 的Gödel数小于第 $i$ 行公式的Gödel数，
  而 $a$ 的Gödel数小于第 $j$ 行公式的Gödel数：
  \begin{multline*}
    \fn{QR}_1(d, i) \defiff \bexists{b < (d)_i}{\bexists{x <
        (d)_i}{\bexists{a < (d)_i}{\bexists{c < (d)_j}{(}}}}
    \\ \fn{Var}(x) \land \fn{Const}(c) \land {} \\
    (d)_i = \Gn{(} \concat
    b \concat \Gn{\lif} \concat \Gn{\lforall} \concat x \concat a
    \concat \Gn{)} \land {}\\
    (d)_j = \Gn{(} \concat b \concat
    \Gn{\lif} \concat \fn{Subst}(a,c,x) \concat \Gn{)} \land {}\\ \fn{Sent}(b)
    \land \fn{Sent}(\fn{Subst}(a,c,x)) \land {} \bforall{k <
      \len{b}}{(b)_k \neq (c)_0})
  \end{multline*}
  这里我们假设 $c$ 和 $x$ 分别是作为项考虑的变元和常项的Gödel数
  （即，不是它们的符号编码）。
  我们通过测试 $\Subst{!A(x)}{c}{x}$ 是否是一个语句来检验 $x$ 是 $!A(x)$ 中唯一的自由变元，
  并通过要求 $!B$ 的每个符号都不同于 $c$ 来确保 $c$ 不在 $!B$ 中出现。

  我们将 \QR{} 的另一版本留作练习。

\item $d$ 是一个正确推导的Gödel数，当且仅当其中的每一行都是公理，或者由肯定前件或 \QR{} 证成。
  因此：
  \[
  \fn{Deriv}(d) \defiff \bforall{i < \len{d}}{(\fn{IsAx}((d)_i) \lor
    \fn{MP}(d,i) \lor \fn{QR}(d, i))}
  \]
\end{enumerate}
\end{proof}

\begin{prob}
如 \olref[inc][art][pax]{prop:followsby} 中那样定义下列关系：
\begin{enumerate}
\item $\fn{IsAx}_{!A \lif (!B \lif (!A \land !B))}(n)$,
\item $\fn{IsAx}_{\lforall[x][!A(x)] \lif !A(t)}(n)$,
\item $\fn{QR}_{2}(d, i)$（针对 \QR{} 的另一版本）。
\end{enumerate}
\end{prob}

\begin{prop}
假设 $\Gamma$ 是一个原始递归的语句集合。
那么关系 $\Prf[\Gamma](x, y)$ 是原始递归的，该关系表达
“$x$ 是从 $\Gamma$ 到 $!A$ 的推导 $\delta$ 的编码，且 $y$ 是 $!A$ 的Gödel数”。
\end{prop}

\begin{proof}
假设“$y \in \Gamma$”由原始递归谓词 $R_\Gamma(y)$ 给出。
我们需要证明关系 $\Prf[\Gamma](x, y)$ 是原始递归的，其中 $\Prf[\Gamma](x, y)$ 成立当且仅当
$y$ 是一个语句 $!A$ 的Gödel数，且 $x$ 是 $!A$ 从 $\Gamma$ 的一个推导的编码。

根据前一命题，性质 $\fn{Deriv}(x)$（它成立当且仅当 $x$ 是一个正确推导 $\delta$ 的编码）是原始递归的。
然而，该定义没有把集合 $\Gamma$ 作为推导中证成一行的额外方式考虑进去。
我们之前对一行是否可由 \QR{} 证成的原始递归检验，也没有顾及常项 $c$ 不允许出现在 $\Gamma$ 中的要求。
要修正我们的定义使其直接考虑 $\Gamma$ 是可能的，但利用 $\fn{Deriv}$ 和演绎定理会更简单。
$\Gamma \Proves !A$ 当且仅当存在某个有穷的语句列表 $!B_1$, \dots, $!B_n \in \Gamma$ 使得 $\{!B_1, \dots, !B_n\} \Proves !A$。
而根据演绎定理，如果 $\Proves (!B_1 \lif (!B_2 \lif
\cdots (!B_n \lif !A)\cdots))$，这就成立。
具有Gödel数 $z$ 的语句是否具有这种形式，是可以原始递归地检验的。
因此，我们不把 $x$ 视为具有Gödel数 $y$ 的语句\emph{从 $\Gamma$} 推导的Gödel数，
而是把 $x$ 视为从 $\emptyset$ 推导出上述形式的一个嵌套条件句的Gödel数。

首先，如果有一个语句序列，我们可以原始递归地构造一个以所有这些语句为前件、以给定语句为后件的条件句：
\begin{align*}
  \fn{hCond}(s, y, 0) & = y \\
  \fn{hCond}(s, y, n+1) & = \Gn{(} \concat (s)_{n} \concat \Gn{\lif}
  \concat \fn{Cond}(s, y, n) \concat \Gn{)}\\
  \fn{Cond}(s, y) & = \fn{hCond}(s, y, \len{s})\\
  \intertext{因此我们可以通过下式定义 $\Prf[\Gamma](x, y)$：}
  \Prf[\Gamma](x, y) & \defiff \bexists{s < \fn{sequenceBound}(x,x)}{(} \\
  &\qquad (x)_{\len{x}-1} = \fn{Cond}(s,y) \land {} \\
  &\qquad\bforall{i<\len{s}}{(s)_i \in \Gamma} \land {} \\
  &\qquad\fn{Deriv}(x)).
\end{align*}
$s$ 的界限是这样得到的：每个 $(s)_i$ 都是该推导最后一行公式的某个子公式的Gödel数，因而小于 $(x)_{\len{x}-1}$。
前件 $!B \in \Gamma$ 的数量，即 $s$ 的长度，小于 $x$ 的最后一行公式的长度。
\end{proof}

\end{document}